\section{Introduction}
\label{sec:introduction}

Atmospheric turbulence couples geometric distortion with spatially varying blur, making single-image restoration ambiguous.
Recovering sharp, undistorted content requires accounting for the physical degradation, particularly when additional frames are unavailable~\cite{mao2022turbnet,jaiswal2023pirn}.

Physics-based simulation connects restoration to image formation.
PiRN~\cite{jaiswal2023pirn} uses differentiable re-degradation for supervision and conditional diffusion for subsequent refinement.
We build on this connection by using the simulator to construct intermediate observations, so that physical image formation also determines how restoration proceeds between stages.

We propose a physics-guided recursive diffusion framework for turbulence restoration.
At each stage, a conditional diffusion restorer based on DiffUIR~\cite{zheng2024diffuir} estimates a clean image, while a separate branch predicts degradation variables.
The simulator renders these estimates at a lower turbulence strength to provide the next stage's input.
During training, rendering at the original strength imposes consistency with the observed image.
Physical strength is conditioned separately from the internal diffusion timestep.
Training progresses from simulated observations to joint optimization of the recursive process, with gradients passing through both sampling and simulation.
Our approach thus combines a simulator-defined restoration path with conditional diffusion and observation-consistency supervision.

\section{Related Work}
\label{sec:related_work}

\paragraph{Atmospheric turbulence restoration.}
Turbulence restoration methods differ in the observations they use and how they handle spatially varying degradation.
For single images, TurbNet~\cite{mao2022turbnet} employs a physics-inspired transformer to estimate distortion and reconstruct scene content.
More recently, D$^2$Turb~\cite{hu2026d2turb} combines depth-aware simulation with separate deblurring and geometric rectification stages, transferring structural features between them.
Video methods can additionally exploit temporal information: DATUM~\cite{zhang2024datum} integrates recurrent aggregation, registration, and lucky-imaging principles, while Turb-Seg-Res~\cite{saha2024turbseg} separates moving foregrounds from static backgrounds before restoration.
These methods highlight the importance of designs tailored to turbulence.
Our setting uses one RGB observation; the intermediate images are generated by the restoration model and simulator, rather than supplied by additional captured frames.

\paragraph{Physics-based simulation and restoration.}
The P2S simulator~\cite{mao2021p2s} maps phase representations to spatial blur representations, enabling efficient synthesis of turbulence-degraded images.
Beyond data generation, differentiable simulation can constrain reconstruction.
PiRN~\cite{jaiswal2023pirn} re-degrades its predicted clean image using the protocol and random realization retained during data synthesis, and penalizes disagreement with the training observation.
In the multi-frame setting, TMFS~\cite{xie2026tmfs} estimates a latent image and degradation parameters using scene-shared correlation functions and per-frame noise maps.
Our formulation also uses a forward model to relate restoration to observation, but predicts the simulator variables from the current restoration state.
The same predictions generate a weaker-turbulence input for the next stage and reconstruct the original-strength observation for supervision.
This use of simulation defines both the recursive transitions and their training constraints; it does not require the estimated variables to uniquely identify the true turbulence realization.

\paragraph{Diffusion for image restoration.}
Diffusion models provide a learned image distribution conditioned on degraded observations.
AT-DDPM~\cite{nair2023atddpm} applies denoising diffusion to turbulence-degraded faces.
For more general scenes, PiRN-SR~\cite{jaiswal2023pirn} conditions a diffusion refiner on the deterministic PiRN reconstruction.
DiffUIR~\cite{zheng2024diffuir} addresses universal restoration through selective hourglass mapping, combining strong image conditioning with a shared distribution across degradation types.
We adopt its conditional residual diffusion formulation as the clean-image estimator within each physical restoration stage.
The simulator then determines the observation supplied to the next stage, while physical strength and diffusion time remain separate conditions.
The distinction from a final diffusion refiner lies in this coupling between repeated conditional restoration and simulator-generated inputs.
